\section{Стандарти оформлення коду та архітектура алгоритмів}

\subsection{Дотримання конвенції Oracle Code Conventions}

У межах виконання лабораторної роботи особливу увагу було приділено стандартам 
інженерного оформлення вихідного коду згідно з вимогами \textit{Oracle Code Conventions}. 
Було опановано та застосовано на практиці такі правила:

\begin{itemize}
    \item \textbf{Угода про іменування (Naming Conventions):} 
    Імена всіх створених компонентів та класів (\texttt{ArrayChecker}, \texttt{FizzBuzzGame}, \texttt{BalanceFinder}) 
    записані у стилі \texttt{UpperCamelCase}. Назви методів та змінних (\texttt{isSortedUp()}, \texttt{totalSum}, \texttt{leftSum}) 
    строго відповідають стилю \texttt{lowerCamelCase}, де перше слово обов'язково є дієсловом для методів та іменником для змінних.
    \item \textbf{Стандартизація коментарів:} Структура кожного файлу містить три рівні документування:
    \begin{enumerate}
        \item \textit{Блок Copyright} у самому верху файлу із зазначенням ліцензії Apache 2.0 та автора.
        \item \textit{JavaDoc коментарі} (\texttt{/** ... */}) перед класами та методами із 
        застосуванням тегів \texttt{@author}, \texttt{@version}, \texttt{@param} та \texttt{@return}.
        \item \textit{Внутрішні коментарі} (\texttt{/* ... */} та \texttt{//}) для 
        пояснення складних ділянок бізнес-логіки з обов'язковим відступом у один порожній рядок перед ними.
    \end{enumerate}
\end{itemize}
